%! ~~~ Packages Setup ~~~ 
\documentclass[]{../../../../tex/classes/styledArticle}
\usepackage{../../../../tex/packages/hebrewSupport}
\usepackage{../../../../tex/packages/mathShortcuts}
\usepackage{../../../../tex/packages/theoremsSupport}
\usepackage{../../../../tex/packages/enfancysections}

%! ~~~ Document ~~~

\author{שחר פרץ}
\title{\textit{לינארית 2א 18 $\sim$ המשפט הספקטרלי להעתקות צמודות לעצמן}}
\date{22 ביוני 2025}
\begin{document}
	\maketitle
	\textbf{מבוסס על הקלטה 17, אלגברה לינארית 2א 2024-2025}
	
	\textbf{מרצה: }ענת אמיר
	
	\theo{אם $V$ ממ''פ ו־$T \co V \to V$ ט''ל צמודה לעצמה, אז הע''ע של $T$ ממשיים. }\begin{proof}
		יהי $0 \neq v \in V$ ו''ע של $T$ שמתאים לע''ע $\lg$. נחשב: 
		\[ \lg v\norm{v}^2 = \mut{\lg v}{v} = \mut{T v}{v} = \mut{v}{Tv} = \mut{v}{\lg v} = \ol \lg \norm{v} \]
		ידוע $v \neq 0$ ולכן $\norm{v} \neq 0$ ונסיק $\lg v = \bar \lg $ ולכן $\lg \in \R$. 
	\end{proof}
	
	\theo{אם $V$ ממ''פ ו־$T \co V \to V$ ט''ל צמודה לעצמה, אז כל זוג $0 \neq u, v \in V$ ע''ע שונים, המתאימים לערכים $\ag, \bg \in \C$, מאונכים זה לזה. }\begin{proof}
		למעשה, מהטענה הקודמת $\ag, \bg \in \R$. כאן $Tu = \ag u, \ Tv = \bg v$, כאשר $\ag = \bg$. נחשב: 
		\[ \ag \mut{u}{v} = \mut{\ag u}{v} = \mut{Tu}{v} = \mut{u}{Tv} = \mut{u}{\bg v} = \bar \bg \mut{u}{v} = \bg \mut{u}{v} \]
		בגלל ש־$\bg \in \R$ מתקיים $\bg = \bar \bg$. ולכן $(\ag - \bg)\mut{v}{u} = 0$ מהעברת אגף וסה''כ $\mut{u}{v} = 0$ ואכן $u \perp v$. 
	\end{proof}
	
	\theo{(המשפט הספקטרלי להעתקה לינארית צמודה לעצמה) יהי $V$ ממ''פ ממימד סופי, ותהי $T \co V \to V$ ט''ל צמודה לעצמה. אז קיים ל־$V$ בסיס אורתוגונלי (או אורתונוגמלי) שמורכב מו''ע של $T$} \begin{proof}
		יהי $m_T(x)$ הפולינום המינימלי של $T$. נציג $m_T(x) = \prod_{i = 1}^{m}(x - \lg_i)^{d_i}$ כאשר $\lg_1 \dots \lg_n$ הע''ע השונים של $T$. מהטענה הקודמת $\lg_1 \dots \lg_n \in \R$. [הערה: התמשתנו במשפט היסודי של האלגברה מעל המרוכבים, והסקנו פירוק מעל $\R$]. בכדי להראות ש־$T$ לכסינה, עלינו להוכיח ש־$\forall 1 \le i \le m \co d_i = 1$. נניח בשלילה שזה לא מתקיים, אזי $m_T(x) = (x - \lg)^2 \cdot p(x)$ כאשר $\lg$ ע''ע כלשהו. כעת, לכל $v \in V$ מתקיים מהיות $T$ צמודה לעצמה (כלומר גם $p(T)$ צמוד לעצמו): 
		\begin{multline*}
			0 = \underbrace{m_T(T)}_{=0}(v) \implies 0 = \mut{m_T(T)(v)}{p(T)(v)} = \mut{(T - \lg I)(p(T)v)}{p(T)v} =\\ \mut{(T - \lg I)(p(T)v)}{(T - \lg I)(p(T)v)} = \norm{(T - \lg I)^2(p(T)v)}^2 = 0
		\end{multline*}
		ולכן $\forall v \in V \co (T - \lg I)(p(T)v) = 0$ ולכן $((x - \lg)(p(x))(T) = 0$ בסתירה למינימליות של $m_T(x)$. נאמר, מכפלת גורמים לינארים שונים, ולכן $T$ לכסינה, ונוכל לפרק את $V$ באמצעות: 
		\[ V = \bigoplus_{i = 1}^{m} \ker (T - \lg_i I) \]
		והמרחכים העצמיים הללו אורתוגונליים זה לזה, מהטענה השנייה שהוכחנו. 
		נבנה בסיס $B_i \subseteq \ker (T - \lg_i)$ וסה''כ $\bigcup_{i = 1}^{m} B_i$ בסיס אורתוגונלי של $T$. 
	\end{proof}
	\theo{יהי $V$ נ''ס מעל $\R$ ותהי $T \co V \to V$ ט''ל. אז $T$ צמודה לעצמה אמ''מ קיים לה בסיס אורתוגונלי מלכסן. }\begin{proof}
		מכיוון אחד, הוכחנו באמצעות המשפט הספקטרלי להעתקות לינאריות צמודות לעצמן. 
		מהכיוון השני, נניח שקיים ל־$V$ בסיס אורתוגונלי מלכסן של ו''ע של $T$. ננרמל לבסיס אורתונורמלי $B = (b_i)_{i = 1}^{n}$ של ו''ע של $T$, המתאימים ל־$\lg_1 \dots \lg_n$. עבור $v, u \in V$, נציג: 
		\[ u = \sum_{i = 1}^{n}\ag_i b_i, \ v = \sum_{i = 1}^{n}\bg_i b_i \]
		ונחשב: 
		\[ \mut{Tu}{v} = \mut{T\cl{\sum_{i = 1}^{m}\ag_i b_i}}{\sum_{i = 1}^{m}\bg_i b_i} = \sum_{i = 1}^{n}\sum_{j = 1}^{n}\ag_i \bg_j \mut{T b_i}{b_j} = \sum_{i = 1}^{n}\sum_{j = 1}^{n}\ag_i \bg_j \lg_i \underbrace{\mut{b_i}{b_j}}_{\dg_{ij}} = \sumni \ag_i \bg_i \lg_i \]
		מהצד השני: 
		\[ \mut{u}{Tv} = \mut{\sumni \ag_i b_i}{T\cl{\sumni \bg_i b_i}} = \sumni \sum_{j = 1}^{n}\ag_i \bg_i \mut{b_i}{Tb_j} = \sum_{i = 1}^{n}\sum_{j = 1}^{n}\ag_i \bg_j \lg_j \underbrace{\mut{b_i}{b_j}}_{\dg_{ij}} = \sum_{i = 1}^{n}\ag_i \bg_j \lg_i \]
		מטרנזטיביות שוויון, הראינו ש־$\mut{Tu}{v} = \mut{u}{Tv}$ ולכן $T$ צמודה לעצמה. השוויון לדלתא של כקוניקר נכונה מאורתוגונליות איברי הבסיס, והבי־לינאריות כי אנחנו מעל הממשיים. המשפט לא נכון מעל מהרוכבים. 
	\end{proof}
	הוכחה שהמשפט לא נכון מעל המרוכבים: ההעתקה $T(x) = ix$ היא העתקה סקלרית לינארית, לכן כל וקטור הוא ו''ע וכל בסיס מלכסן, בסיס אורתונורמלי כלשהו יהיה בסיס מלכסן על אף שההעתקה לא צמודה לעצמה, אלא אנטי־הרמיטית. 
	
	מכאן ואילך המרצה מוכיחה את המשפט הספקטרלי ללא המשפט היסודי של האלגברה. לשם כך, צריך להראות שהפולינום המינימלי מתפצל למכפלה של גורמים לינארים מעל המרוכבים. 
	
	\theo{אם $p(x)$ פולינום אי־שלילי, אז נוכל להציגו כ־$p(x) = \sum_{i = 1}^{k}g_i(x)^2 + c$, כאשר $c \ge 0$, וכמו כן $c > 0$ אמ''מ $p(x)$ פולינום חיובי. }\begin{proof}
		נבחין ש־$p(x)$ בהכרח ממעלה זוגית מהיותו אי־שלילי, כי לפולינום אי־זוגי מתקיים שהגבולות באינסוף מחליפים סימן. נוכיח באינדוקציה 
		
		טוב כאן נאמס לי, אני עובר להקלטה הבאה. 
		TODO: להשלים את ההוכחה הזו. 
	\end{proof}
	
	\ndoc
\end{document}
